% IMMUTABLE RESEARCH REPORT::HASHLIT1::2026-10-07_234214
\documentclass[11pt]{article}
\usepackage[T1]{fontenc}
\usepackage[margin=1in]{geometry}
\usepackage{amsmath,amssymb,amsthm,booktabs,array,longtable,hyperref}
\setlength{\emergencystretch}{3em}
\hypersetup{colorlinks=true,linkcolor=blue,citecolor=blue,urlcolor=blue}

\newtheorem{theorem}{Theorem}
\newtheorem{lemma}[theorem]{Lemma}
\newtheorem{corollary}[theorem]{Corollary}
\newtheorem{proposition}[theorem]{Proposition}
\newtheorem{definition}[theorem]{Definition}
\newcommand{\SAT}{\operatorname{SAT}}
\newcommand{\poly}{\operatorname{poly}}
\newcommand{\mc}{\operatorname{mc}}
\newcommand{\Var}{\operatorname{Var}}

\title{Exact Fingerprints, Hash-Consing, and SAT State Caching\\
Primary-Source Transfers and Limits for the WRAPPER Program}
\author{RESEARCH agent report HASHLIT1}
\date{7 October 2026}

\begin{document}
\maketitle

\begin{abstract}
This report asks whether deterministic fingerprints, perfect hashing,
adaptive fingerprints, canonical DAG memoization, or SAT component caches can
turn a short WRAPPER label into an exact all-input merge rule.  The answer is
model-dependent.  A fixed $d$-digit label used as the sole merge certificate
is exact only when at most $10^d$ inequivalent states are reachable.  A short
label can instead serve as a bucket selector, but then correctness comes from
an exact structural comparison or a source-checkable semantic certificate,
whose acquisition, replay, time, and workspace must be charged.  Primary
literature supplies strong exact mechanisms for already acquired finite key
sets, ordered decision diagrams, bounded-width compilation, and
source-provenance verification on already compiled decision-DNNFs.  It does
not supply a general-input SAT compressor or a fixed-width exact semantic
fingerprint.  Therefore choosing four, six, or eight output digits changes
constants in collision rates and bucket sizes, but cannot by itself change a
polynomial or exponential resource exponent.  The missing general WRAPPER
contraction remains \textbf{OPEN}; consequently this report does not establish
$P\ne NP$.
\end{abstract}

\section{Status, target, and quantifiers}

The target considered here is an exact deterministic merge procedure for a
family of WRAPPER states.  At input scale $N$, let $\mathcal S_N$ be the set
of states that the declared construction can reach, and let $\equiv_N$ be the
continuation equivalence that licenses all future reuse.  Equality of the
current SAT answer is generally weaker than continuation equivalence: two
states can both be satisfiable and still behave differently after later
restrictions or recursive calls.

The load-bearing desired statement has the quantifier form
\[
  \exists\,\text{one uniform exact constructor }M\quad
  \forall N\quad\forall x\in\{0,1\}^N,
\]
followed by simultaneous bounds on generated states, encoded child size,
local work, exact verification, recursive calls, and peak workspace.  A
finite experiment, an average-case hashing theorem, or a structure theorem
for bounded width does not establish this all-input statement.

The concrete question is whether a label
\[
 h_N:\mathcal S_N\longrightarrow [10^d]
\]
with $d\in\{4,6,8\}$, or an adaptive extension of such a label, can supply
the missing exact contraction.  The distinction
\[
 \text{existence}\ne\text{acquisition}\ne
 \text{verification}\ne\text{decision cost}
\]
is maintained throughout.

\section{The exact fixed-range obstruction}

\begin{lemma}[Fixed-range exact-merge obstruction]\label{lem:fixed-range}
Suppose a deterministic wrapper merges two reachable states whenever their
$d$-digit labels agree, without any further exact verification.  If every
merge is sound for $\equiv_N$, then $h_N$ induces an injection
\[
 \mathcal S_N/\!\equiv_N\;\hookrightarrow\;[10^d].
\]
Consequently
\[
  10^d\ge \left|\mathcal S_N/\!\equiv_N\right|,
  \qquad
  d\ge
  \left\lceil\log_{10}\left|\mathcal S_N/\!\equiv_N\right|\right\rceil.
  \tag{1}
\]
\end{lemma}

\begin{proof}
If two distinct $\equiv_N$-classes had the same label, the wrapper would
perform an unsound merge.  Selecting one representative of each equivalence
class therefore gives distinct labels.  The conclusion follows by the
pigeonhole principle.
\end{proof}

\begin{corollary}[Bucket-selector alternative]\label{cor:selector}
If equal labels only nominate candidates and a separate verifier $V_N(x,y)$
must accept before a merge, the wrapper may remain exact with a noninjective
label.  In that case exactness comes from $V_N$, not from the $d$ digits, and
the number of verifier calls, source access, certificate length, replay time,
and verifier workspace are part of the resource bound.
\end{corollary}

The obstruction is independent of how the digits are produced.  It applies
to decimal truncation, modular reduction, a fixed cryptographic digest,
universal hashing after the coins are fixed, or a truncated product of prime
labels.

\subsection{Four, six, and eight digits}

The exact range and information content are as follows.

\begin{center}
\begin{tabular}{@{}rrrr@{}}
\toprule
Decimal digits $d$ & Range $10^d$ & $\log_2(10^d)$ & Stored binary bits \\
\midrule
4 & $10^4$ & $13.2877\ldots$ & 14 \\
6 & $10^6$ & $19.9316\ldots$ & 20 \\
8 & $10^8$ & $26.5754\ldots$ & 27 \\
\bottomrule
\end{tabular}
\end{center}

All three widths are $\Theta(1)$ as $N\to\infty$.  If the number of
reachable inequivalent states is $E(N)=N^{a+o(1)}$, exact merge-on-label
requires
\[
 d(N)\ge (a+o(1))\log_{10}N.
 \tag{2}
\]
If $E(N)=2^{N^{b+o(1)}}$, it requires
\[
 d(N)\ge (\log_{10}2)N^{b+o(1)}.
 \tag{3}
\]
Thus fixed digits do not scale to either family.

The same conclusion appears directly in exponent arithmetic.  If $Q(N)$
candidate states incur balanced all-pairs bucket work on the scale
\[
 \frac{Q(N)^2}{10^d},
\]
then for $Q(N)=N^{a+o(1)}$ and fixed $d$ the exponent remains $2a$.
To lower it by $b>0$ through range size alone requires
\[
 10^{d(N)}=N^{b+o(1)},
 \qquad d(N)=(b+o(1))\log_{10}N.
 \tag{4}
\]
Likewise
\[
 \frac{2^k}{10^d}=\Theta(2^k)
\]
for fixed $d$.  A strict exponential-rate change needs $d=\Omega(k)$.

\section{Deterministic equality and streaming fingerprints}

The fundamental exact lower bound is elementary and is also the standard
communication-complexity explanation for why random fingerprints help only
when error is allowed.

\begin{lemma}[Exact streaming equality]\label{lem:eq-stream}
Let a deterministic one-pass streaming machine receive $x\#y$, where
$x,y\in\{0,1\}^n$, and accept exactly when $x=y$.  Immediately after the
separator, the machine has at least $2^n$ possible states.  Hence it uses at
least $n$ bits of state.
\end{lemma}

\begin{proof}
If distinct $x,x'$ led to the same state after their first halves, append
$y=x$.  Determinism would force identical behavior on $x\#x$ and
$x'\#x$, although exactly one is an equality instance.
\end{proof}

This is the one-way form of deterministic communication complexity for
$n$-bit equality.  Lalonde, Mande, and de Wolf record deterministic cost $n$
and analyze the sharp small-error randomized and quantum regimes
\cite{lalonde2023}.  Babu et al. give a model-matched multipass transfer:
their Theorem~2 proves that every deterministic $p$-pass streaming
recognizer for $1$-turn-$\mathrm{Dyck}_2$ uses $\Omega(n/p)$ space, by
communicating the streaming memory in a reduction from equality
\cite{babu2011}.

\paragraph{Transfer to the wrapper.}
If a wrapper must compare two arbitrary $n$-bit exact state descriptions
from their fingerprints alone, Lemma~\ref{lem:eq-stream} forbids an exact
constant-size deterministic fingerprint.  This is a theorem transfer only
for the declared equality task.  It is not a lower bound on all SAT
algorithms and does not show that the wrapper must materialize arbitrary
state descriptions.

\section{Universal hashing and randomized equality}

Carter and Wegman define a universal class of functions from a universe into
a finite range so that for every fixed pair $x\ne y$, a uniformly selected
function collides on that pair with probability at most the reciprocal of
the range size \cite{carterwegman1979}.  In the notation here,
\[
 \Pr_h[h(x)=h(y)]\le 10^{-d}.
 \tag{5}
\]
This is a probability over the selected function for a fixed pair.  It is
not a deterministic statement about every input after the function is
published, and the original paper explicitly distinguishes the case in
which the user can exploit knowledge of the chosen function.

For $C$ fixed comparisons, the union bound gives
\[
 \Pr[\text{at least one collision}]\le C10^{-d}.
 \tag{6}
\]
To make this at most $\delta$ requires
\[
 d\ge \left\lceil\log_{10}(C/\delta)\right\rceil.
 \tag{7}
\]
If $C=N^q$ and $\delta=N^{-c}$, the right side is
$(q+c)\log_{10}N+O(1)$.  If $L$ live histories may all be compared, the
ordinary all-pairs estimate is
\[
 \Pr[\text{some pair collides}]
 \le \binom{L}{2}10^{-d}.
 \tag{8}
\]
For $L=2^k$, deterministic collision freedom requires at least
$\lceil k\log_{10}2\rceil$ digits, while the all-pairs randomized bound needs
roughly twice that many plus the desired error term.

Lalonde, Mande, and de Wolf show that public-coin equality can use about
$\log_2(1/\varepsilon)$ communication bits, whereas the private-coin regime
also pays for establishing the randomness and depends on $n$
\cite{lalonde2023}.  Thus a public shared seed, a private streaming seed, and
a fixed deterministic hash are different models.  A fixed seed leaves a
deterministic map with permanent colliding pairs.

\paragraph{Status.}
Universal hashing is a \textbf{KNOWN THEOREM} for per-pair randomized
collision bounds.  Importing it as an exact all-input merge certificate is a
\textbf{FAILED TRANSFER}.  It can support a bounded-error wrapper only if the
final theorem permits randomness and the seed, adaptivity, total comparison
count, and effective error are all charged.

\section{Las Vegas collision verification}

Karp and Rabin's string-matching algorithms illustrate the sound use of a
short fingerprint in an exact computation \cite{karprabin1987}.  Their
always-correct variants verify each fingerprint hit against the original
strings, and one variant chooses a new fingerprint after a false match.
Their Theorem~5 gives expected $O(n+m)$ running time for the stated pattern
and text lengths.  Under their modulus choice, Theorem~6 bounds repeated
false matches after reinitialization.

The relevant transfer is:
\[
 \text{fingerprint hit}
 \longrightarrow
 \text{candidate}
 \longrightarrow
 \text{exact source comparison}.
 \tag{9}
\]
Correctness survives collisions because the original objects remain
accessible.  The expected bound depends on the paper's random selection
model.  Worst-case exact source replay can still be large.  Therefore this
mechanism is a useful wrapper design pattern, but it is not lossless
constant-width compression.

\section{Perfect hashing after the key set is known}

Fredman, Koml\'os, and Szemer\'edi construct a static dictionary for a known
finite set of $L$ keys using $O(L)$ space and $O(1)$ worst-case membership
queries; the standard FKS construction uses expected linear randomized
construction time \cite{fks1984}.  Dynamic perfect hashing retains linear
space and $O(1)$ worst-case lookup with expected amortized $O(1)$ updates
\cite{dietzfelbinger1994}.

These results are exact because collisions among the stored keys are
eliminated or resolved while the original keys are available.  They do not
give a fixed-width injective map for arbitrarily many keys.  An index among
$L$ table positions requires $\Theta(\log L)$ bits, and the construction must
read or enumerate the set it perfect-hashes.

Hagerup, Miltersen, and Pagh give a deterministic counterpart
\cite{hagerup2001}.  Their Theorem~1.1, on a unit-cost word RAM with
multiplication and word length $w$, constructs a static dictionary for $n$
$w$-bit keys in $O(n\log n)$ time and $O(n)$ words, with constant lookup.
The algorithm is weakly nonuniform: it uses word-size constants depending on
$w$ that are not known to be computable in $w^{O(1)}$ time.  Their universe
reduction maps the known key set injectively into $O(\log n)$ bits, not into a
constant number of bits.

Miltersen gives a particularly clean circuit statement
\cite{miltersen1997}: for every known set
$S\subseteq\{0,1\}^w$ of size $n$, there is a Boolean circuit of size
$O(w\log w)$, with $w$ inputs and $2\log n$ output bits up to ordinary
rounding, whose function is one-to-one on $S$.  This proves that an exact
deterministic perfect fingerprint can exist for an already specified finite
live set.  It also displays the two limitations needed here: output width
grows as $\Theta(\log n)$, and the theorem does not cheaply discover an
unknown exponentially large state family or semantic quotient.

\paragraph{Wrapper consequence.}
Perfect hashing may compact references to states that the wrapper has already
acquired.  It does not establish that few states are generated, that semantic
equivalence is cheap to decide, or that acquisition and storage fit the
desired exponent.

\section{Adaptive and extendible fingerprints}

Adaptive approximate-membership structures respond to discovered false
positives by changing or extending fingerprints.  Bender et al. prove both a
lower bound and a matching-style construction in their model
\cite{bender2018}.  For $n$ keys from a sufficiently large universe and a
constant sustained false-positive probability, their Theorem~5 requires
\[
 \Omega\!\left(\min\{n\log n,\,n\log\log u\}\right)
\]
total bits with high probability.  Their broom filter in Theorem~6 uses
approximately $n\log(1/\varepsilon)+O(n)$ local bits with high probability,
supports constant-time operations with high probability, and appends bits to
separate discovered collisions.  The complete keys remain available in a
remote representation, and false positives can trigger remote updates.

Naor and Yogev study a stronger adversarial model in which future queries
depend on previous answers \cite{naoryogev2019}.  For computationally bounded
adversaries, strong resilience is connected to cryptographic assumptions
such as one-way functions or pseudorandom permutations.  For an unbounded
adversary making $t$ adaptive queries, their construction uses
\[
 O\!\left(n\log(1/\varepsilon)+t\right)
\]
bits.  Ordinary per-fixed-query collision estimates do not automatically
survive adaptive reuse.

These are approximate randomized data structures.  They show how useful
systems grow information, retain full keys elsewhere, or assume
cryptographic pseudorandomness.  None turns four, six, or eight digits into
an exact deterministic identifier for an unbounded wrapper state family.

\section{Exact structural hash-consing and canonical DAGs}

Bryant's reduced ordered Boolean decision diagrams are the canonical exact
hash-consing example \cite{bryant1986}.  For a fixed variable ordering,
Theorem~1 states that every Boolean function has a unique reduced function
graph up to isomorphism, and every other ordered function graph for that
function has more vertices.  Bottom-up reduction identifies a node by the
exact tuple
\[
  (\text{variable},\text{low-child ID},\text{high-child ID}).
 \tag{10}
\]
Nodes with identical tuples are merged.  An implementation may hash the tuple
to find a bucket, but it compares the exact tuple to resolve collisions.

If $Q$ unique nodes remain, their identifiers require $\Theta(\log Q)$ bits.
Moreover, Bryant explicitly allows worst-case exponential graph size; many
practical functions are small, but canonicity does not imply polynomial
succinctness for every Boolean function.

\paragraph{Exact transfer.}
Bryant's rule transfers exactly to a wrapper whose residual states possess a
proved canonical recursive representation with cheaply computed child IDs.
The missing theorem is precisely that every general SAT input produces only
polynomially many such canonical states and that constructing them fits the
target resource ledger.  Without that theorem, OBDD hash-consing is a
restricted representation result rather than a general SAT contraction.

\section{Information-theoretically succinct exact dictionaries}

Raman, Raman, and Satti store an $n$-element subset of an $m$-element
universe in
\[
 \mathcal B(n,m)+o(n)+O(\log\log m)
 \quad\text{bits},\qquad
 \mathcal B(n,m)=\left\lceil\log_2\binom{m}{n}\right\rceil,
 \tag{11}
\]
while supporting constant-time rank and select; in the cell-probe model the
additive term can be removed \cite{raman2007}.  The leading term is the
information-theoretic minimum.

This theorem shows that exact storage can approach optimal information
content.  It does not eliminate that content.  A wrapper may encode its
already known state set succinctly, but arbitrary state families cannot be
represented by a constant number of digits unless the family itself has a
uniformly bounded number of possibilities or distinguishing information is
stored elsewhere.

\section{SAT component caching and compiled search traces}

Bacchus, Dalmao, and Pitassi analyze DPLL with exact component caching for
\#SAT \cite{bacchus2009}.  For an $n$-variable instance of branchwidth $w$,
their full-cache algorithm has a deterministic execution in
\[
 n^{O(1)}2^{O(w)}\quad\text{time},
 \tag{12}
\]
while their linear-space variant uses space linear in the instance size and
time
\[
 2^{O(w\log n)}.
 \tag{13}
\]
The cache keys are exact residual formula components; hash values, if used in
an implementation, do not authorize merging nonidentical components.

This supplies a rigorous time--space route on bounded-width families.  For
$w=O(1)$ it is polynomial; for $w=\Theta(n)$ it gives no general polynomial
bound.  It therefore supports an achievable-region theorem parameterized by
branchwidth, not an unrestricted SAT compression theorem.

Huang and Darwiche identify the propositional languages generated by several
forms of exhaustive DPLL search \cite{huangdarwiche2007}.  Depending on
variable-ordering, decomposition, and caching restrictions, traces correspond
to FBDD, OBDD, or a decision-DNNF subclass.  This correspondence allows known
succinctness lower bounds for a target representation to become limitations
on the corresponding search-trace algorithm.  It does not lower-bound every
SAT algorithm, because the transfer is tied to the declared trace model.

\section{Source-backed semantic hashing on decision-DNNF}

Biere, Lagniez, and Lonca give the closest current literature analogue to the
proposed forward wrapper \cite{biere2026}.  The input is not an arbitrary CNF
alone.  It consists of a decision-DNNF $\Psi$ already compiled from and
representing an original CNF $\Sigma$, with provenance linking circuit nodes
to the source.

For each node $u$, they compute the semantic signature
\[
 H(u)=\langle\Var(u),\mc(u)\rangle,
 \tag{14}
\]
where $\mc(u)$ is the exact model count.  Equal signatures are necessary but
not sufficient for equivalence, so they only create candidate buckets.
Their Proposition~3 states that, for two nodes $N_1,N_2$ of the compiled
circuit, the two entailments
\[
 \operatorname{circ}(N_1)\models\operatorname{cnf}(N_2,\Sigma),
 \qquad
 \operatorname{circ}(N_2)\models\operatorname{cnf}(N_1,\Sigma)
 \tag{15}
\]
imply circuit equivalence.  Clause entailment is polynomial-time on DNNF, so
the source projections provide an exact polynomial verifier in this compiled
model.  Proposition~4 states that the resulting circuit has no two distinct
nodes computing equivalent Boolean functions.

This is an exact theorem transfer for the pattern
\[
 \boxed{\text{semantic filter} + \text{source provenance}
        + \text{exact mutual-entailment verification}.}
 \tag{16}
\]
It also states the principal limitations.  A general CNF may require an
exponentially large decision-DNNF before this minimization begins.  Exact
model counts have growing bit length.  Buckets may still contain many
nonequivalent nodes and trigger many polynomial-time comparisons.  The paper
notes that equivalence of two arbitrary decision-DNNFs remains open.  The
provenance theorem therefore does not provide a general CNF compiler or a
fixed-digit exact hash.

Truncating $\mc(u)$ to four, six, or eight decimal digits leaves the
verification sound, because equation~(15) still controls merging.  It can
only enlarge candidate buckets and increase verification work.  It cannot
improve the worst-case exponent without a separate bucket-size theorem.

\section{Canonical TDDs and bounded-treewidth compilation}

Capelli et al. introduce Tree Decision Diagrams (TDDs), a canonical
generalization of OBDDs \cite{capelli2026}.  Three theorems are directly
relevant.

\begin{description}
\item[Theorem 11.] Given a TDD $C$ of width $k$ over variables $X$, a
minimal canonical representation $m(C)$ can be computed in
$\poly(k)|X|$ time.

\item[Theorem 16.] A CNF with $n$ variables, $m$ clauses, and primal or
incidence treewidth $k$ can be compiled into a TDD of width at most $2^k$ in
$2^{O(k)}mn$ time.

\item[Theorem 18.] A Boolean circuit $C$ of treewidth $k$ can be compiled
into a TDD of width $2^{O(k)}$ in time
$2^{O(k)}|X||C|$.
\end{description}

The minimization is exact and deterministic: duplicate subfunctions become
``twins'' that can be contracted in the canonical representation.  This is a
genuine positive route for bounded treewidth or, more generally, bounded
factor width.  The same paper proves exponential lower bounds for TDD size
for some functions (Theorem~21).  It therefore cannot be promoted to a
polynomial-size representation of all SAT instances.

\section{Algebraic fingerprints and deterministic limits}

Schwartz's polynomial identity test and the related Schwartz--Zippel lemma
show that a nonzero degree-$D$ polynomial evaluated at a uniformly random
point from a finite set $S$ vanishes with probability at most $D/|S|$
\cite{schwartz1980}.  This is another randomized fingerprint: a fixed field
or fixed number of output digits does not give zero error as the degree,
number of candidates, or number of tests grows.

Kabanets and Impagliazzo show that broad deterministic polynomial-time
derandomization of polynomial identity testing would imply major arithmetic
or Boolean circuit lower bounds \cite{kabanetsimpagliazzo2004}.  This result
must not be overapplied: equality of explicit strings has simple exact
methods, and restricted canonical representations may be tractable.  Its
proper use here is only to reject an unsupported claim that every circuit
state admits a cheap deterministic algebraic fingerprint with all the power
of randomized identity testing.

\section{Compressed fingerprints and verification cost}

Bille et al. study fingerprints inside grammar-compressed strings
\cite{bille2017}.  They obtain small-space Monte Carlo fingerprint query
structures, while their always-correct Las Vegas construction pays a
potentially much larger preprocessing cost to certify collision freedom.
For a string of uncompressed length $N$ and grammar size $n$, their stated
Las Vegas preprocessing bound includes
$O(N^2/(n\log N))$ in the general setting.

This is not a SAT lower bound.  It is a model-matched warning: when a short
fingerprint indexes a compressed object, certifying that the fingerprint is
collision-free may depend on the uncompressed source scale.  A wrapper must
therefore charge provenance generation and collision certification rather
than count only the stored digest.

\section{Prime-product labels}

Suppose the four truth-status labels
\[
 \mathsf{TRUE},\quad\mathsf{NOT\ TRUE},\quad
 \mathsf{FALSE},\quad\mathsf{NOT\ FALSE}
\]
receive four fixed primes and a history is represented by their product.
Unique factorization then recovers only the four multiplicities.  It does not
recover which variable received which status or the order and recursion
depth at which statuses appeared.  Distinct histories with the same four
counts collide even before decimal truncation.

Assigning a distinct prime to every pair
$(\text{variable or position},\text{status})$ restores injectivity for a
fixed ordered history.  For $n$ positions, the relevant primes are among the
first $4n$, each has $O(\log n)$ bits, and a product of $n$ selected primes
has $\Theta(n\log n)$ bits in the worst case.  Recursive or order-sensitive
histories also require position- or depth-specific factors.  Reducing the
product modulo $10^4$, $10^6$, or $10^8$ returns to
Lemma~\ref{lem:fixed-range}.

The prime scheme is therefore a source-checkable exact encoding when its full
integer and factor assignment are retained.  It is not constant-space
compression.

\section{Transfer classification}

\begin{longtable}{@{}p{0.22\textwidth}p{0.23\textwidth}p{0.47\textwidth}@{}}
\toprule
Mechanism & Classification & Exact wrapper consequence \\
\midrule
\endhead
Deterministic equality lower bound & THEOREM TRANSFER & Exact comparison of arbitrary
$n$-bit states from a one-pass summary needs $n$ bits.  This applies to the
declared equality subtask, not every SAT algorithm. \\
\addlinespace
Universal hashing & FAILED TRANSFER for exactness & Provides per-fixed-pair
random collision bounds.  Exact all-input merging still needs verification
or a growing range. \\
\addlinespace
Karp--Rabin verification & THEOREM TRANSFER & A short label may safely nominate
candidates if the original source is retained and every match is verified.
Expected bounds require the paper's randomness model. \\
\addlinespace
Perfect hashing & THEOREM TRANSFER after acquisition & Gives exact fast lookup
for an already known finite key set.  It does not discover or shrink the
state family, and IDs grow as $\log L$. \\
\addlinespace
Adaptive filters & ANALOGY / CANDIDATE & Show how systems append bits and use
remote full keys.  They are approximate and randomized, so they do not
certify exact semantic merging. \\
\addlinespace
Bryant OBDD reduction & THEOREM TRANSFER under fixed order & Exact tuple
hash-consing gives a canonical DAG.  General size can be exponential. \\
\addlinespace
Component caching & THEOREM TRANSFER for branchwidth $w$ & Gives exact
$n^{O(1)}2^{O(w)}$ full-cache time or linear-space
$2^{O(w\log n)}$ time.  No unrestricted polynomial bound follows. \\
\addlinespace
Decision-DNNF semantic hashing & THEOREM TRANSFER after compilation & Source
provenance plus mutual entailment gives exact merges.  Compilation and bucket
cost may be exponential. \\
\addlinespace
Canonical TDD minimization & THEOREM TRANSFER for bounded structural width &
Provides deterministic canonical contraction and FPT compilation on bounded
treewidth; some functions still require exponential TDDs. \\
\addlinespace
Fixed four/six/eight digits & FALSIFIED/RETIRED as sole exact certificate &
The range is constant, so it cannot distinguish a growing family of
inequivalent states or alter an asymptotic exponent. \\
\bottomrule
\end{longtable}

\section{A source-checkable exact merge schema}

The literature supports the following conditional schema.  It is a precise
interface, not a completed general-SAT construction.

\begin{definition}[Certified bucket merge]
For every generated state $s$, the forward wrapper produces:
\[
  \bigl(B(s),\,R(s),\,P(s)\bigr),
\]
where $B(s)$ is a cheap bucket label, $R(s)$ is a replayable exact
representation or source reference, and $P(s)$ is provenance sufficient for
an exact verifier.  States in different buckets are kept distinct.  States
in one bucket are merged only when a deterministic verifier
\[
 V(R(s),P(s);R(t),P(t))=1
\]
proves continuation equivalence.
\end{definition}

\begin{proposition}[Conditional correctness]
If $V$ is sound for continuation equivalence and every reuse depends only on
the equivalence class certified by $V$, every accepted merge is sound.  If
$V$ is also complete on the generated state family, the retained states are
exactly one representative per equivalence class.
\end{proposition}

\begin{proof}
Soundness prevents an accepted pair from crossing equivalence classes.
Completeness ensures that every pair in one class can eventually be joined.
The bucket label affects search order only.
\end{proof}

To become the missing WRAPPER theorem, this schema still needs all-input
bounds for the quantities below.  None is supplied by the digit choice.

\begin{description}
\item[Generation.] Total number $Q(N)$ of states and the cost of producing
$B,R,P$ for each state.

\item[Encoded size.] Maximum encoded size of every child passed to recursion
or to a SAT procedure.

\item[Bucket burden.] A deterministic bound on
\[
  C_{\rm cand}(N)=\sum_z \binom{|B_z|}{2}
\]
or on the actual candidate schedule.

\item[Verification.] Worst-case time and workspace of each $V$ call, including
source regeneration, exact comparison, entailment, or certificate checking.

\item[Storage/replay.] Peak retained representation and provenance space, or
the charged cost of regenerating evicted data.

\item[Recursion.] Total branching, recursive query size, recombination cost,
and closure of the same bounds at every smaller scale.

\item[Randomness, if any.] Seed length, adaptive-query model, total failure
probability over the entire recursive computation, and whether the endpoint
lower-bound theorem admits Monte Carlo or Las Vegas algorithms.
\end{description}

\section{What this report changes}

The earlier repository red-team report already established the internal
fixed-hash pigeonhole and bucket-replay countermodels.  This report does not
claim a new lower bound against all SAT algorithms.  It adds a primary-source
map showing which pieces are known theorems and exactly where their models
stop.

\begin{description}
\item[Adopted.] Fixed-range injectivity obstruction; exact hashes as bucket
indices; source replay must be charged; bounded-width caching and canonical
representations as restricted positive routes.

\item[Strengthened.] The exact deterministic obstruction is now tied to
communication/streaming equality; deterministic and randomized perfect
hashing are separated from semantic discovery; the 2026 decision-DNNF and TDD
results give concrete source-backed and canonical mechanisms.

\item[Rejected.] The claim that four, six, or eight digits can alone produce
an exact asymptotic contraction; the claim that universal hashing remains an
exact guarantee after fixing the seed; the claim that perfect hashing avoids
enumerating or storing the live state family.

\item[Still open.] A uniform general-CNF constructor that cheaply generates
sufficiently few smaller instances, supplies exact continuation-equivalence
certificates, and meets the required joint time--space exponent on every
input.
\end{description}

No root-main statement can be promoted from \textbf{OPEN} on the basis of
this literature.  The strongest defensible positive theorem is restricted:
source-backed exact merging is available after decision-DNNF compilation,
and canonical deterministic compilation is fixed-parameter tractable for
bounded-treewidth CNFs and circuits.

\section{Next concrete target}

The next useful theorem is not a longer fixed digest.  It is an all-input
bound for a particular certified bucket merge:
\begin{center}
\fbox{\begin{minipage}{0.84\textwidth}
Give a uniform source-checkable rule that, before any current-scale SAT
query, produces smaller encoded children; prove deterministic bounds on
the number of children, candidate pairs, exact verification/replay work,
retained provenance, recursive query size, and peak workspace.
\end{minipage}}
\end{center}
The decision-DNNF provenance rule and the TDD twin-contraction rule are the
best current templates.  The missing bridge is to obtain their benefits
without first paying exponential general-input compilation cost.

\section*{Primary-source URLs}

\begin{description}
\item[Equality and streaming.] \url{https://doi.org/10.4230/LIPIcs.FSTTCS.2023.32};
\url{https://arxiv.org/abs/1104.0848}.

\item[Universal and perfect hashing.] 
\url{https://doi.org/10.1016/0022-0000(79)90044-8};
\url{https://doi.org/10.1145/828.1884};
\url{https://doi.org/10.1137/S0097539791194094};
\url{https://doi.org/10.1006/jagm.2001.1171};
\url{https://doi.org/10.7146/brics.v4i17.18813}.

\item[Fingerprints and adaptivity.] 
\url{https://doi.org/10.1147/rd.312.0249};
\url{https://arxiv.org/abs/1711.01616};
\url{https://arxiv.org/abs/1412.8356}.

\item[Canonical DAGs and succinct dictionaries.]
\url{https://doi.org/10.1109/TC.1986.1676819};
\url{https://arxiv.org/abs/0705.0552}.

\item[SAT caching and compilation.]
\url{https://eccc.weizmann.ac.il/report/2003/003/};
\url{https://arxiv.org/abs/1401.3458};
\url{https://arxiv.org/abs/1110.2734};
\url{https://doi.org/10.24963/ijcai.2026/420};
\url{https://doi.org/10.4230/LIPIcs.SAT.2026.10}.

\item[Algebraic and compressed fingerprints.]
\url{https://doi.org/10.1145/322217.322225};
\url{https://doi.org/10.1007/s00037-004-0182-6};
\url{https://arxiv.org/abs/1305.2777}.
\end{description}

\begin{thebibliography}{99}

\bibitem{babu2011}
A.~Babu, N.~Limaye, J.~Radhakrishnan, and G.~Varma.
\newblock Streaming algorithms for language recognition problems.
\newblock \emph{Theoretical Computer Science}, 494:13--23, 2013.
\newblock Preprint: \url{https://arxiv.org/abs/1104.0848}.

\bibitem{bacchus2009}
F.~Bacchus, S.~Dalmao, and T.~Pitassi.
\newblock Solving \#SAT and Bayesian inference with backtracking search.
\newblock \emph{Journal of Artificial Intelligence Research}, 34:391--442,
2009.
\newblock \url{https://arxiv.org/abs/1401.3458}.

\bibitem{bender2018}
M.~A. Bender, M.~Farach-Colton, M.~Goswami, R.~Johnson, S.~McCauley, and
S.~Singh.
\newblock Bloom filters, adaptivity, and the dictionary problem.
\newblock In \emph{FOCS 2018}.
\newblock \url{https://arxiv.org/abs/1711.01616}.

\bibitem{biere2026}
A.~Biere, J.-M. Lagniez, and E.~Lonca.
\newblock Efficient minimization of decision-DNNF circuits via semantic
hashing and provenance tracking.
\newblock In \emph{IJCAI 2026}, pages 3774--3782.
\newblock \url{https://doi.org/10.24963/ijcai.2026/420}.

\bibitem{bille2017}
P.~Bille, P.~H. Cording, I.~L. G\o rtz, B.~Sach, H.~W. Vildh\o j, and
S.~Vind.
\newblock Fingerprints in compressed strings.
\newblock \emph{Journal of Computer and System Sciences}, 86:171--180, 2017.
\newblock \url{https://arxiv.org/abs/1305.2777}.

\bibitem{bryant1986}
R.~E. Bryant.
\newblock Graph-based algorithms for Boolean function manipulation.
\newblock \emph{IEEE Transactions on Computers}, C-35(8):677--691, 1986.
\newblock \url{https://doi.org/10.1109/TC.1986.1676819}.

\bibitem{capelli2026}
F.~Capelli, Y.~Choi, S.~Mengel, M.~Mu\~noz, and G.~Van den Broeck.
\newblock A canonical generalization of OBDD.
\newblock In \emph{SAT 2026}, LIPIcs 377, Article 10.
\newblock \url{https://doi.org/10.4230/LIPIcs.SAT.2026.10}.

\bibitem{carterwegman1979}
J.~L. Carter and M.~N. Wegman.
\newblock Universal classes of hash functions.
\newblock \emph{Journal of Computer and System Sciences}, 18(2):143--154,
1979.
\newblock \url{https://doi.org/10.1016/0022-0000(79)90044-8}.

\bibitem{dietzfelbinger1994}
M.~Dietzfelbinger, A.~R. Karlin, K.~Mehlhorn, F.~Meyer auf der Heide,
H.~Rohnert, and R.~E. Tarjan.
\newblock Dynamic perfect hashing: upper and lower bounds.
\newblock \emph{SIAM Journal on Computing}, 23(4):738--761, 1994.
\newblock \url{https://doi.org/10.1137/S0097539791194094}.

\bibitem{fks1984}
M.~L. Fredman, J.~Koml\'os, and E.~Szemer\'edi.
\newblock Storing a sparse table with $O(1)$ worst case access time.
\newblock \emph{Journal of the ACM}, 31(3):538--544, 1984.
\newblock \url{https://doi.org/10.1145/828.1884}.

\bibitem{hagerup2001}
T.~Hagerup, P.~B. Miltersen, and R.~Pagh.
\newblock Deterministic dictionaries.
\newblock \emph{Journal of Algorithms}, 41(1):69--85, 2001.
\newblock \url{https://doi.org/10.1006/jagm.2001.1171}.

\bibitem{huangdarwiche2007}
J.~Huang and A.~Darwiche.
\newblock The language of search.
\newblock \emph{Journal of Artificial Intelligence Research}, 29:191--219,
2007.
\newblock \url{https://doi.org/10.1613/JAIR.2097}.

\bibitem{kabanetsimpagliazzo2004}
V.~Kabanets and R.~Impagliazzo.
\newblock Derandomizing polynomial identity tests means proving circuit lower
bounds.
\newblock \emph{Computational Complexity}, 13:1--46, 2004.
\newblock \url{https://doi.org/10.1007/s00037-004-0182-6}.

\bibitem{karprabin1987}
R.~M. Karp and M.~O. Rabin.
\newblock Efficient randomized pattern-matching algorithms.
\newblock \emph{IBM Journal of Research and Development}, 31(2):249--260,
1987.
\newblock \url{https://doi.org/10.1147/rd.312.0249}.

\bibitem{lalonde2023}
O.~Lalonde, N.~Mande, and R.~de Wolf.
\newblock Tight bounds for the randomized and quantum communication
complexities of equality with small error.
\newblock In \emph{FSTTCS 2023}, LIPIcs 284, Article 32.
\newblock \url{https://doi.org/10.4230/LIPIcs.FSTTCS.2023.32}.

\bibitem{miltersen1997}
P.~B. Miltersen.
\newblock Error correcting codes, perfect hashing circuits, and deterministic
dynamic dictionaries.
\newblock \emph{BRICS Report Series}, 4(17), 1997.
\newblock \url{https://doi.org/10.7146/brics.v4i17.18813}.

\bibitem{naoryogev2019}
M.~Naor and E.~Yogev.
\newblock Bloom filters in adversarial environments.
\newblock \emph{ACM Transactions on Algorithms}, 15(3), Article 34, 2019.
\newblock \url{https://doi.org/10.1145/3306193}.

\bibitem{raman2007}
R.~Raman, V.~Raman, and S.~R. Satti.
\newblock Succinct indexable dictionaries with applications to encoding
$k$-ary trees, prefix sums and multisets.
\newblock \emph{ACM Transactions on Algorithms}, 3(4), Article 43, 2007.
\newblock \url{https://arxiv.org/abs/0705.0552}.

\bibitem{schwartz1980}
J.~T. Schwartz.
\newblock Fast probabilistic algorithms for verification of polynomial
identities.
\newblock \emph{Journal of the ACM}, 27(4):701--717, 1980.
\newblock \url{https://doi.org/10.1145/322217.322225}.

\end{thebibliography}

\paragraph{Report verification status.}
Literature claims above are labeled only as transfers within their stated
models.  No general SAT contraction or complexity-class separation is
promoted.  Compilation status is recorded outside this immutable scientific
body by the creating run.

% END_OF_VERIFIED_REPORT::RESEARCH::HASHLIT1
\end{document}
